% Generated mechanically from frozen input inputs/p10/ja/spaces-more-cohomology.tex.
% Do not edit this generated file; edit the bound locale source instead.

% OK, start here.
%

\setcounter{chapter}{83}
\chapter{空間のコホモロジー詳論}


\phantomsection
\label{spaces-more-cohomology-section-phantom}





\section{序論}
\label{spaces-more-cohomology-section-introduction}

\noindent
% P10-E0147
本章では、『空間のコホモロジー』節
\ref{spaces-cohomology-section-introduction} で始めた議論を続ける。
また本章は、スキームに対するエタール・コホモロジーの章、すなわち
『エタール・コホモロジー』節
\ref{etale-cohomology-section-introduction} の代数空間版とみなすこともできる。

\medskip\noindent
実際、本章の主な目的は、スキームについて既に証明された結果を
代数空間の言葉に移し替えることである。本章の結果の一部は
\cite{Kn} にも見られる。





\section{規約}
\label{spaces-more-cohomology-section-conventions}

\noindent
常に、すべてのスキームが大 fppf サイト $\Sch_{fppf}$ に
含まれるものと仮定する。また、考えるすべての環 $A$ は、
$\Spec(A)$ がこの大サイトの対象に（同型で）なるという性質をもつものとする。

\medskip\noindent
$S$ をスキームとし、$X$ を $S$ 上の代数空間とする。
本章および次章では、$X$ とそれ自身との積（$S$ 上の代数空間の圏における積）を
$X \times X$ ではなく $X \times_S X$ と書く。







\section{スキームの場合の結果の移送}
\label{spaces-more-cohomology-section-api}

\noindent
本節では、スキームについての結果から、（表現可能な）代数空間および
代数空間の（表現可能な）射についての結果がどのように従うかを簡潔に説明する。
準連接加群については、さらに強いことが成り立つ
（スキーム上の準連接加群のエタール・コホモロジーは
Zariski コホモロジーと一致するためである）。これは既に
『空間のコホモロジー』節
\ref{spaces-cohomology-section-higher-direct-image} で論じた。

\medskip\noindent
$S$ をスキームとし、$X$ を $S$ 上の代数空間とする。
いま $X$ がスキーム $X_0$ によって表現されると仮定する
（不格好だが一時的な記法であり、通常は単に「$X$ はスキームである」と言う）。
このとき $X$ と $X_0$ の小エタールサイトは同じである：
$$
X_\etale = (X_0)_\etale
$$
これは『代数空間の性質』節
\ref{spaces-properties-section-etale-site} で指摘されている。
さらに、$f : X \to Y$ が $S$ 上の表現可能な代数空間の射であり、
$f_0 : X_0 \to Y_0$ が $f$ を表現するスキームの射ならば、
誘導される小エタール・トポスの射は一致する：
$$
\xymatrix{
\Sh(X_\etale) \ar[rr]_{f_{small}} \ar@{=}[d] & &
\Sh(Y_\etale) \ar@{=}[d] \\
\Sh((X_0)_\etale) \ar[rr]^{(f_0)_{small}} & &
\Sh((Y_0)_\etale)
}
$$
『代数空間の性質』補題
\ref{spaces-properties-lemma-functoriality-etale-site} および
『位相』補題 \ref{topologies-lemma-morphism-big-small-etale} を参照されたい。

\medskip\noindent
したがって、スキームのエタール・コホモロジーと、それに対応する
代数空間のエタール・コホモロジーとの間にはまったく差がない。
スキームの射に沿う高次順像についても同様である。
実際、$f : X \to Y$ が $S$ 上の代数空間の射であり、
（スキームで）表現可能ならば、$X_\etale$ 上の層 $\mathcal{F}$ の
高次順像 $R^if_*\mathcal{F}$ は $Y$ 上エタール局所的に計算できる
（『サイト上のコホモロジー』補題
\ref{sites-cohomology-lemma-higher-direct-images}）。
したがって、計算や証明はしばしば $Y$ と $X$ がスキームである場合に帰着する。

\medskip\noindent
本章では、以上を以後断りなく用いる。他の位相についても同じことが成り立つので、
ここではコホモロジーに関する形で明示的に補題として述べておく。

\begin{lemma}
\label{spaces-more-cohomology-lemma-compare-cohomology-other-topologies}
$S$ をスキームとする。
% P10-E0148
$\tau \in \{\etale, fppf, ph\}$ とする（ここにさらに加える）。包含関手
$$
(\Sch/S)_\tau \longrightarrow (\textit{Spaces}/S)_\tau
$$
は特殊余連続関手であり
（『サイト』定義 \ref{sites-definition-special-cocontinuous-functor}）、
したがって両トポスを同一視する。
\end{lemma}

\begin{proof}
この関手は完全忠実であり、また任意の代数空間はスキームによる
エタール被覆をもつので、『サイト』補題
\ref{sites-lemma-equivalence} の条件は直ちに確認できる。
\end{proof}







\section{固有基底変換}
\label{spaces-more-cohomology-section-proper-base-change}

\noindent
代数空間に対する固有基底変換定理は、少し作業すれば、
スキームに対する固有基底変換定理と Chow の補題から従う。

\begin{lemma}
\label{spaces-more-cohomology-lemma-surjective-proper}
$S$ をスキームとする。$f : Y \to X$ を $S$ 上の代数空間の
全射な固有射とする。$\mathcal{F}$ を $X_\etale$ 上の層とする。
このとき $\mathcal{F} \to f_*f^{-1}\mathcal{F}$ は単射であり、その像は
二つの写像 $f_*f^{-1}\mathcal{F} \to g_*g^{-1}\mathcal{F}$ の等化子である。
ここで $g$ は構造射 $g : Y \times_X Y \to X$ である。
\end{lemma}

\begin{proof}
$S$ 上の代数空間の任意の全射 $f : Y \to X$ に対し、写像
$\mathcal{F} \to f_*f^{-1}\mathcal{F}$ は単射である。
実際、$X$ の幾何点 $\overline{x}$ を取るとき、$\overline{x}$ の上にある
$Y$ の幾何点 $\overline{y}$ を選び、
$$
\mathcal{F}_{\overline{x}} \to
(f_*f^{-1}\mathcal{F})_{\overline{x}} \to
(f^{-1}\mathcal{F})_{\overline{y}} = \mathcal{F}_{\overline{x}}
$$
を考えればよい。最後の等号については『代数空間の性質』補題
\ref{spaces-properties-lemma-stalk-pullback} を参照されたい。

\medskip\noindent
第二の主張は $X$ 上エタール局所的なので、
% P10-E0605
$Y$ がアフィンスキームであると仮定してよいし、実際そう仮定する。

\medskip\noindent
$Z$ がスキームであるような全射固有射 $Z \to Y$ を選ぶ。
『空間のコホモロジー』補題
\ref{spaces-cohomology-lemma-weak-chow} を参照されたい。
$Z \to X$ に対する結果から $Y \to X$ に対する結果が従う。
$Z \to X$ はスキームの全射固有射であり、したがって ph 被覆である
（『位相』補題 \ref{topologies-lemma-surjective-proper-ph}）から、
$Z \to X$ に対する結果は『エタール・コホモロジー』補題
\ref{etale-cohomology-lemma-describe-pullback-pi-ph} から従う
（実際、ある意味でこの補題と同値である）。
\end{proof}

\begin{lemma}
\label{spaces-more-cohomology-lemma-h0-proper-over-henselian-pair}
$(A, I)$ を Hensel 対とする。$X$ を $A$ 上の代数空間とし、
構造射 $f : X \to \Spec(A)$ が固有であるとする。
$i : X_0 \to X$ を $X \times_{\Spec(A)} \Spec(A/I)$ の包含とする。
$X_\etale$ 上の任意の層 $\mathcal{F}$ に対し
$\Gamma(X, \mathcal{F}) = \Gamma(X_0, i^{-1}\mathcal{F})$ が成り立つ。
\end{lemma}

\begin{proof}
$Y$ がスキームであるような全射固有射 $Y \to X$ を選ぶ。
『空間のコホモロジー』補題
\ref{spaces-cohomology-lemma-weak-chow} を参照されたい。図式
$$
\xymatrix{
\Gamma(X_0, \mathcal{F}_0) \ar[r] &
\Gamma(Y_0, \mathcal{G}_0) \ar@<1ex>[r] \ar@<-1ex>[r] &
\Gamma((Y \times_X Y)_0, \mathcal{H}_0) \\
\Gamma(X, \mathcal{F}) \ar[r] \ar[u] &
\Gamma(Y, \mathcal{G}) \ar@<1ex>[r] \ar@<-1ex>[r] \ar[u] &
\Gamma(Y \times_X Y, \mathcal{H}) \ar[u]
}
$$
を考える。
% P10-E0149
ここで $\mathcal{G}$ および $\mathcal{H}$ は、それぞれ $\mathcal{F}$ の
$Y$ および $Y \times_X Y$ への逆像であり、添字 $0$ は
$\Spec(A/I)$ への基底変換を表す。スキームの場合
（『エタール・コホモロジー』補題
\ref{etale-cohomology-lemma-h0-proper-over-henselian-pair}）により、
中央と右の垂直射は全単射である。補題 \ref{spaces-more-cohomology-lemma-surjective-proper} により、
左の垂直射も全単射であることが従う。
\end{proof}

\begin{lemma}
\label{spaces-more-cohomology-lemma-h0-proper-over-henselian-local}
$A$ を Hensel 局所環とする。$X$ を $A$ 上の代数空間とし、
$f : X \to \Spec(A)$ が固有射であるとする。
$X_0 \subset X$ を閉点上の $f$ のファイバーとする。
$X_\etale$ 上の任意の層 $\mathcal{F}$ に対し
$\Gamma(X, \mathcal{F}) = \Gamma(X_0, \mathcal{F}|_{X_0})$ が成り立つ。
\end{lemma}

\begin{proof}
これは補題 \ref{spaces-more-cohomology-lemma-h0-proper-over-henselian-pair} の特別な場合である。
\end{proof}

\begin{lemma}
\label{spaces-more-cohomology-lemma-proper-base-change-f-star}
$S$ をスキームとする。
% P10-E0150
$f : X \to Y$ および $g : Y' \to Y$ を $S$ 上の代数空間の射とし、
$f$ は固有であると仮定する。
$X' = Y' \times_Y X$ とおき、射影を $f' : X' \to Y'$ および
$g' : X' \to X$ とする。$\mathcal{F}$ を $X_\etale$ 上の任意の層とする。
このとき
$g^{-1}f_*\mathcal{F} = f'_*(g')^{-1}\mathcal{F}$ である。
\end{lemma}

\begin{proof}
問題は $Y'$ 上エタール局所的である。
スキーム $V$ と全射エタール射 $V \to Y$ を選ぶ。
さらにスキーム $V'$ と全射エタール射 $V' \to V \times_Y Y'$ を選ぶ。
すると $Y'$ を $V'$ で、$Y$ を $V$ で置き換えてよい。
したがって $Y$ と $Y'$ はスキームであると仮定できる。
さらに $Y$ と $Y'$ 上 Zariski 局所的に議論してよいので、
$Y$ と $Y'$ はアフィンスキームであると仮定できる。

\medskip\noindent
$Y$ と $Y'$ がアフィンスキームであると仮定する。
$X_1$ がスキームであるような全射固有射 $h_1 : X_1 \to X$ を選ぶ。
『空間のコホモロジー』補題
\ref{spaces-cohomology-lemma-weak-chow} を参照されたい。
$X_2 = X_1 \times_X X_1$ とおき、構造射を $h_2 : X_2 \to X$ と記す。
これはスキームであることに注意する。スキームの場合
（『エタール・コホモロジー』補題
\ref{etale-cohomology-lemma-proper-base-change-f-star}）により、
補題は Cartesian 図式
$$
\vcenter{
\xymatrix{
X'_1 \ar[r] \ar[d] & X_1 \ar[d] \\
Y' \ar[r] & Y
}
}
\quad\text{および}\quad
\vcenter{
\xymatrix{
X'_2 \ar[r] \ar[d] & X_2 \ar[d] \\
Y' \ar[r] & Y
}
}
$$
と層 $\mathcal{F}_i = (X_i \to X)^{-1}\mathcal{F}$ について成り立つ。
補題 \ref{spaces-more-cohomology-lemma-surjective-proper} により完全列
$0 \to \mathcal{F} \to h_{1, *}\mathcal{F}_1 \to h_{2, *}\mathcal{F}_2$
を得る。また $X'_2 = X'_1 \times_{X'} X'_1$ なので、
$(g')^{-1}\mathcal{F}$ についても同様である。したがって
% P10-E0151
補題が成り立つことが分かる（若干の詳細は省略する）。
\end{proof}

\noindent
$S$ をスキームとする。$f : Y \to X$ を $S$ 上の代数空間の射とする。
% P10-E0152
$\overline{x} : \Spec(k) \to S$ を幾何点とする。
$\overline{x}$ における $f$ のファイバーとは、$\Spec(k)$ 上の代数空間
$Y_{\overline{x}} = \Spec(k) \times_{\overline{x}, X} Y$ のことである。
$\mathcal{F}$ が $Y_\etale$ 上の層ならば、
$\mathcal{F}_{\overline{x}} = p^{-1}\mathcal{F}$ を、
$\mathcal{F}$ の $(Y_{\overline{x}})_\etale$ への逆像と記す。
ここで $p : Y_{\overline{x}} \to Y$ は射影である。
以下では集合 $\Gamma(Y_{\overline{x}}, \mathcal{F}_{\overline{x}})$ を考える。

\begin{lemma}
\label{spaces-more-cohomology-lemma-proper-pushforward-stalk}
$S$ をスキームとする。$f : Y \to X$ を $S$ 上の代数空間の固有射とし、
$\overline{x} \to X$ を幾何点とする。
$Y_\etale$ 上の任意の層 $\mathcal{F}$ に対し、標準写像
$$
(f_*\mathcal{F})_{\overline{x}} \longrightarrow
\Gamma(Y_{\overline{x}}, \mathcal{F}_{\overline{x}})
$$
は全単射である。
\end{lemma}

\begin{proof}
これは補題 \ref{spaces-more-cohomology-lemma-proper-base-change-f-star} の特別な場合である。
\end{proof}

\begin{theorem}
\label{spaces-more-cohomology-theorem-proper-base-change}
$S$ をスキームとし、
$$
\xymatrix{
X' \ar[r]_{g'} \ar[d]_{f'} & X \ar[d]^f \\
Y' \ar[r]^g & Y
}
$$
を $S$ 上の代数空間の Cartesian 正方形とする。$f$ は固有であると仮定する。
$\mathcal{F}$ を $X_\etale$ 上のアーベル捩れ層とする。このとき基底変換写像
$$
g^{-1}Rf_*\mathcal{F} \longrightarrow Rf'_*(g')^{-1}\mathcal{F}
$$
は同型である。
\end{theorem}

\begin{proof}
この証明では、スキームに対する固有基底変換定理の証明で用いた議論の一部を
繰り返す。詳しくは『エタール・コホモロジー』節
\ref{etale-cohomology-section-proper-base-change} を参照されたい。

\medskip\noindent
主張は $Y'$ および $Y$ 上エタール局所的なので、$Y$ と $Y'$ はともに
アフィンスキームであると仮定してよい。特に、これで $f$ が表現可能な場合の
定理も証明されたことに注意する（以下でこれを用いる）。

\medskip\noindent
各 $n \geq 1$ に対し、$\mathcal{F}[n]$ を $n$ で零化される
$\mathcal{F}$ の切断の部分層とする。このとき
$\mathcal{F} = \colim \mathcal{F}[n]$ である。
『空間のコホモロジー』補題
\ref{spaces-cohomology-lemma-colimit-cohomology} により、関手
$g^{-1}R^pf_*$ および $R^pf'_*(g')^{-1}$ はフィルター付き余極限と可換である。
したがって、$\mathcal{F}$ が $n$ で零化される場合を証明すれば十分である。

\medskip\noindent
$\mathcal{F} \to \mathcal{I}^\bullet$ を
$\mathbf{Z}/n\mathbf{Z}$-加群の入射層による分解とする。
補題 \ref{spaces-more-cohomology-lemma-proper-base-change-f-star} により
$g^{-1}f_*\mathcal{I}^\bullet = f'_*(g')^{-1}\mathcal{I}^\bullet$
であることに注意する。Leray の非輪状性補題
（『導来圏』補題 \ref{derived-lemma-leray-acyclicity}）を適用すると、
$p > 0$ および $m \in \mathbf{Z}$ に対し
$R^pf'_*(g')^{-1}\mathcal{I}^m = 0$ を証明すれば十分である。

\medskip\noindent
$Z$ がスキームであるような全射固有射 $h : Z \to X$ を選ぶ。
『空間のコホモロジー』補題
\ref{spaces-cohomology-lemma-weak-chow} を参照されたい。
$h^{-1}\mathcal{I}^m \to \mathcal{J}$ を単射とし、$\mathcal{J}$ を
$Z_\etale$ 上の $\mathbf{Z}/n\mathbf{Z}$-加群の入射層とする。
$h$ は全射なので、写像 $\mathcal{I}^m \to h_*\mathcal{J}$ は単射である
（補題 \ref{spaces-more-cohomology-lemma-surjective-proper} を参照）。
$\mathcal{I}^m$ は入射的なので、$\mathcal{I}^m$ は
$h_*\mathcal{J}$ の直和因子である。したがって、望む消滅を
$h_*\mathcal{J}$ について証明すれば十分である。

\medskip\noindent
% P10-E0153
$h'$ を $g$ による基底変換とし、射影を $g'' : Z' \to Z$ と記す。
スペクトル系列
$$
E_2^{p, q} = R^pf'_* R^qh'_* (g'')^{-1}\mathcal{J}
$$
があり、$R^{p + q}(f' \circ h')_*(g'')^{-1}\mathcal{J}$ に収束する。
$h$ および $f \circ h$ は（スキームで）表現可能なので、
求める結果がこれらについて成り立つことは既に分かっている。
したがって、このスペクトル系列において $q > 0$ なら
$E_2^{p, q} = 0$ であり、$p + q > 0$ なら
$R^{p + q}(f' \circ h')_*(g'')^{-1}\mathcal{J} = 0$ である。
ゆえに $p > 0$ なら $E_2^{p, 0} = 0$ である。さて、補題
\ref{spaces-more-cohomology-lemma-proper-base-change-f-star} により
$$
E_2^{p, 0} = R^pf'_* h'_* (g'')^{-1}\mathcal{J} =
R^pf'_* (g')^{-1}h_*\mathcal{J}
$$
である。これで証明が完了する。
\end{proof}

\begin{lemma}
\label{spaces-more-cohomology-lemma-proper-base-change}
$S$ をスキームとし、
$$
\xymatrix{
X' \ar[r]_{g'} \ar[d]_{f'} & X \ar[d]^f \\
Y' \ar[r]^g & Y
}
$$
を $S$ 上の代数空間の Cartesian 正方形とする。$f$ は固有であると仮定する。
$E \in D^+(X_\etale)$ のコホモロジー層が捩れ層であるとする。
このとき基底変換写像 $g^{-1}Rf_*E \to Rf'_*(g')^{-1}E$ は同型である。
\end{lemma}

\begin{proof}
これは、スペクトル系列
$$
E_2^{p, q} = R^pf_*H^q(E)
\quad\text{および}\quad
{E'}_2^{p, q} = R^pf'_*(g')^{-1}H^q(E)
$$
を用いた固有基底変換定理（定理 \ref{spaces-more-cohomology-theorem-proper-base-change}）の
簡単な帰結である。これらはそれぞれ $R^nf_*E$ および
$R^nf'_*(g')^{-1}E$ に収束する。スペクトル系列は『導来圏』補題
\ref{derived-lemma-two-ss-complex-functor} で構成されている。
若干の詳細は省略する。
\end{proof}

\begin{lemma}
\label{spaces-more-cohomology-lemma-proper-base-change-stalk}
$S$ をスキームとする。$f : X \to Y$ を代数空間の固有射とし、
$\overline{y} \to Y$ を幾何点とする。
\begin{enumerate}
\item $X_\etale$ 上のアーベル捩れ層 $\mathcal{F}$ に対し
$(R^nf_*\mathcal{F})_{\overline{y}} =
H^n_\etale(X_{\overline{y}}, \mathcal{F}_{\overline{y}})$ である。
\item コホモロジー層が捩れ層である $E \in D^+(X_\etale)$ に対し\\
$(R^nf_*E)_{\overline{y}} = H^n_\etale(X_{\overline{y}}, E_{\overline{y}})$
である。
\end{enumerate}
\end{lemma}

\begin{proof}
主張において、$\mathcal{F}_{\overline{y}}$ は $\mathcal{F}$ の
$X_{\overline{y}} = \overline{y} \times_Y X$ への逆像を表す。
$\overline{y} \to Y$ による逆像は $\mathcal{F}$ の茎を与えるので、
第一の主張は定理 \ref{spaces-more-cohomology-theorem-proper-base-change} の特別な場合である。
第二の主張は補題 \ref{spaces-more-cohomology-lemma-proper-base-change} の特別な場合である。
\end{proof}

\begin{lemma}
\label{spaces-more-cohomology-lemma-base-change-separably-closed}
$k'/k$ を分離閉体の拡大とする。$X$ を $k$ 上の固有代数空間とし、
$\mathcal{F}$ を $X$ 上のアーベル捩れ層とする。このとき $q \geq 0$ に対し、
$H^q_\etale(X, \mathcal{F}) \to
H^q_\etale(X_{k'}, \mathcal{F}|_{X_{k'}})$ は同型である。
\end{lemma}

\begin{proof}
これは定理 \ref{spaces-more-cohomology-theorem-proper-base-change} の特別な場合である。
\end{proof}







\section{大トポスと小トポスの比較}
\label{spaces-more-cohomology-section-compare}

\noindent
$S$ をスキーム、$X$ を $S$ 上の代数空間とする。
『空間上のトポロジー』補題
\ref{spaces-topologies-lemma-at-the-bottom-etale} では、比較射
$\pi_X : (\textit{Spaces}/X)_\etale \to X_{spaces, \etale}$ および
$i_X : \Sh(X_\etale) \to \Sh((\textit{Spaces}/X)_\etale)$
を導入した。これらはトポスの射として
$\pi_X \circ i_X = \text{id}$ を満たし、$\pi_{X, *} = i_X^{-1}$ である。
より一般に、$f : Y \to X$ が $(\textit{Spaces}/X)_\etale$ の対象ならば、
$f_{small} = \pi_X \circ i_f$ を満たす射
$i_f : \Sh(Y_\etale) \to \Sh((\textit{Spaces}/X)_\etale)$
が存在する。『空間上のトポロジー』補題
\ref{spaces-topologies-lemma-put-in-T-etale} および
\ref{spaces-topologies-lemma-morphism-big-small-etale} を参照されたい。
『空間上のトポロジー』注意
\ref{spaces-topologies-remark-change-topologies-ringed} では、これらを環付きサイトの射
$$
\pi_X :
((\textit{Spaces}/X)_\etale, \mathcal{O})
\to
(X_{spaces, \etale}, \mathcal{O}_X)
$$
ならびに環付きトポスの射
$$
i_X :
(\Sh(X_\etale), \mathcal{O}_X)
\to
(\Sh((\textit{Spaces}/X)_\etale), \mathcal{O})
$$
および
$$
i_f :
(\Sh(Y_\etale), \mathcal{O}_Y)
\to
(\Sh((\textit{Spaces}/X)_\etale, \mathcal{O}))
$$
へ拡張した。制限 $i_X^{-1} = \pi_{X, *}$（『トポロジー』定義
\ref{topologies-definition-restriction-small-etale} を参照）は
$\mathcal{O}$ を $\mathcal{O}_X$ に移すことに注意する。
同様に、$i_f^{-1}$ は $\mathcal{O}$ を $\mathcal{O}_Y$ に移す。
『空間上のトポロジー』注意
\ref{spaces-topologies-remark-change-topologies-ringed} を参照されたい。
したがって、$(\textit{Spaces}/X)_\etale$ 上の任意の
$\mathcal{O}$-加群 $\mathcal{F}$ に対し
$i_X^*\mathcal{F} = i_X^{-1}\mathcal{F}$ および
$i_f^*\mathcal{F} = i_f^{-1}\mathcal{F}$ である。特に $i_X^*$ と $i_f^*$ は
完全関手である。関手 $i_X^*$ はしばしば
$\mathcal{F} \mapsto \mathcal{F}|_{X_\etale}$ と書かれる
（これは『空間上のトポロジー』定義
\ref{spaces-topologies-definition-restriction-small-etale} の記法と矛盾しない）。

\begin{lemma}
\label{spaces-more-cohomology-lemma-describe-pullback}
$S$ をスキーム、$X$ を $S$ 上の代数空間とする。
$\mathcal{F}$ を $X_\etale$ 上の層とする。このとき
$\pi_X^{-1}\mathcal{F}$ は、$(\textit{Spaces}/X)_\etale$ の
$f : Y \to X$ に対する規則
$$
(\pi_X^{-1}\mathcal{F})(Y) = \Gamma(Y_\etale, f_{small}^{-1}\mathcal{F})
$$
で与えられる。さらに、
% P10-E0154
$\pi_Y^{-1}\mathcal{F}$ は平滑被覆、シントミック被覆、fppf 被覆、
fpqc 被覆、および ph 被覆に関する層条件を満たす。
\end{lemma}

\begin{proof}
逆像の推移性と $f_{small} = \pi_X \circ i_f$（上記参照）から
$i_f^{-1} \pi_X^{-1}\mathcal{F} = f_{small}^{-1}\mathcal{F}$ を得る。
これにより、$\pi_X^{-1}$ が補題に記した記述をもつことが分かる。

\medskip\noindent
$\pi_X^{-1}\mathcal{F}$ が ph トポロジーに対する層であることを示すには、
『空間上のトポロジー』補題
\ref{spaces-topologies-lemma-characterize-sheaf} により、
$X$ 上の代数空間の全射固有射 $V \to U$ に対し
$(\pi_X^{-1}\mathcal{F})(U)$ が二つの写像
$(\pi_X^{-1}\mathcal{F})(V) \to (\pi_X^{-1}\mathcal{F})(V \times_U V)$
の等化子であることを示せば十分である。これは補題
\ref{spaces-more-cohomology-lemma-surjective-proper} で示した。

\medskip\noindent
平滑被覆、シントミック被覆、および fppf 被覆の場合は、
『空間上のトポロジー』補題
\ref{spaces-topologies-lemma-zariski-etale-smooth-syntomic-fppf-ph}
により ph 被覆の場合から従う。

\medskip\noindent
$\mathcal{U} = \{U_i \to U\}_{i \in I}$ を $X$ 上の代数空間の
fpqc 被覆とする。$s_i \in (\pi_X^{-1}\mathcal{F})(U_i)$ を、
$U_i \times_U U_j$ 上で一致する切断とする。$U_i$ 上で $s_i$ に
制限される一意な $s \in (\pi_X^{-1}\mathcal{F})(U)$ が存在することを
証明しなければならない。場合 I：$U$ と $U_i$ がスキームである場合。
これは『エタール・コホモロジー』補題
\ref{etale-cohomology-lemma-describe-pullback} から従う。
場合 II：$U$ がスキームである場合。ここで、$T_i$ がスキームであるような
全射エタール射 $T_i \to U_i$ を選ぶ。このとき
$\mathcal{T} = \{T_i \to U\}$ はスキームによる fpqc 被覆であり、
場合 I により結果は $\mathcal{T}$ に対して成り立つ。
これが $\mathcal{U}$ に対する結果を導くことの確認は省略する。
場合 III：一般の場合。$W$ がスキームであるような全射エタール射
$W \to U$ をとる。このとき
$\mathcal{W} = \{U_i \times_U W \to W\}$ は、スキーム $W$ の
（代数空間による）fpqc 被覆である。場合 II により、
% P10-E0155
結果は $\mathcal{W}$ に対して成り立つ。
これが $\mathcal{U}$ に対する結果を導くことの確認は省略する。
\end{proof}

\begin{lemma}
\label{spaces-more-cohomology-lemma-compare-injectives}
$S$ をスキームとし、$Y \to X$ を $(\textit{Spaces}/S)_\etale$ の射とする。
\begin{enumerate}
\item $\mathcal{I}$ が $\textit{Ab}((\textit{Spaces}/X)_\etale)$ において
入射的ならば、
\begin{enumerate}
\item $i_f^{-1}\mathcal{I}$ は $\textit{Ab}(Y_\etale)$ において入射的であり、
\item $\mathcal{I}|_{X_\etale}$ は $\textit{Ab}(X_\etale)$ において入射的である。
\end{enumerate}
\item $\mathcal{I}^\bullet$ が
$\textit{Ab}((\textit{Spaces}/X)_\etale)$ における K-入射複体ならば、
\begin{enumerate}
\item $i_f^{-1}\mathcal{I}^\bullet$ は $\textit{Ab}(Y_\etale)$ における
K-入射複体であり、
\item $\mathcal{I}^\bullet|_{X_\etale}$ は $\textit{Ab}(X_\etale)$ における
K-入射複体である。
\end{enumerate}
\end{enumerate}
加群に対する対応する主張は成り立たない。
\end{lemma}

\begin{proof}
(1)(b) と (2)(b) は、制限関手 $\pi_{X, *} = i_X^{-1}$ が
完全関手 $\pi_X^{-1}$ の右随伴であることから形式的に従う。
『ホモロジー』補題 \ref{homology-lemma-adjoint-preserve-injectives} および
『導来圏』補題 \ref{derived-lemma-adjoint-preserve-K-injectives} を参照されたい。

\medskip\noindent
(1)(a) と (2)(a) は二通りに示せる。第一の証明：$i_f^{-1}$ が
完全関手 $i_{f, !}$ の右随伴であることを用いる。この関手は集合の層について
『トポロジー』補題 \ref{topologies-lemma-put-in-T-etale} で、
アーベル層について『サイト上の加群』補題
\ref{sites-modules-lemma-g-shriek-adjoint} で構成される。
『サイト上の加群』補題
\ref{sites-modules-lemma-exactness-lower-shriek} において、これが完全であることが示される。
第二の証明：『トポロジー』補題
\ref{topologies-lemma-morphism-big-small-etale} で示される
$i_f = i_Y \circ f_{big}$ を用いる。$f_{big}$ は局所化なので、
これによる逆像は入射対象と K-入射対象を保つ。
『サイト上のコホモロジー』補題
\ref{sites-cohomology-lemma-cohomology-of-open} および
\ref{sites-cohomology-lemma-restrict-K-injective-to-open} を参照されたい。
次に、既に証明した (1)(b) と (2)(b) を関手 $i_Y^{-1}$ に適用して結論を得る。

\medskip\noindent
加群の場合の反例については、『エタール・コホモロジー』補題
\ref{etale-cohomology-lemma-compare-injectives} を参照されたい。
\end{proof}

\noindent
$S$ をスキームとし、$f : Y \to X$ を $S$ 上の代数空間の射とする。
『空間上のトポロジー』補題
\ref{spaces-topologies-lemma-morphism-big-small-etale} (3) の可換図式から、
環付きサイトの可換図式
$$
\xymatrix{
(Y_{spaces, \etale}, \mathcal{O}_Y) \ar[d]_{f_{spaces, \etale}} &
((\textit{Spaces}/Y)_\etale, \mathcal{O}) \ar[d]^{f_{big}} \ar[l]^{\pi_Y} \\
(X_{spaces, \etale}, \mathcal{O}_X) &
((\textit{Spaces}/X)_\etale, \mathcal{O}) \ar[l]_{\pi_X}
}
$$
を得る。これは $f_{small}^\sharp$、$f_{big}^\sharp$、
$\pi_X^\sharp$、および $\pi_Y^\sharp$ の定義を書き下せば容易に分かる。
特に、$(\textit{Spaces}/Y)_\etale$ 上の任意の層 $\mathcal{F}$ に対し
\begin{equation}
\label{spaces-more-cohomology-equation-compare-big-small}
(f_{big, *}\mathcal{F})|_{X_\etale} =
f_{small, *}(\mathcal{F}|_{Y_\etale})
\end{equation}
であり、$\mathcal{F}$ が $\mathcal{O}$-加群の層ならば、
(\ref{spaces-more-cohomology-equation-compare-big-small}) は $X_\etale$ 上の
$\mathcal{O}_X$-加群の同型である。

\begin{lemma}
\label{spaces-more-cohomology-lemma-compare-higher-direct-image}
$S$ をスキーム、$f : Y \to X$ を $S$ 上の代数空間の射とする。
\begin{enumerate}
\item $D((\textit{Spaces}/Y)_\etale)$ の $K$ に対し、\\
$
(Rf_{big, *}K)|_{X_\etale} = Rf_{small, *}(K|_{Y_\etale})
$
が $D(X_\etale)$ において成り立つ。
\item $D((\textit{Spaces}/Y)_\etale, \mathcal{O})$ の $K$ に対し、\\
$
(Rf_{big, *}K)|_{X_\etale} = Rf_{small, *}(K|_{Y_\etale})
$
が $D(\textit{Mod}(X_\etale, \mathcal{O}_X))$ において成り立つ。
\end{enumerate}
より一般に、$g : X' \to X$ を $(\textit{Spaces}/X)_\etale$ の対象とし、
ファイバー積
$$
\xymatrix{
Y' \ar[r]_{g'} \ar[d]_{f'} & Y \ar[d]^f \\
X' \ar[r]^g & X
}
$$
を考える。このとき
\begin{enumerate}
\item[(3)] $D((\textit{Spaces}/Y)_\etale)$ の $K$ に対し、
$i_g^{-1}(Rf_{big, *}K) = Rf'_{small, *}(i_{g'}^{-1}K)$
が $D(X'_\etale)$ において成り立つ。
\item[(4)] $D((\textit{Spaces}/Y)_\etale, \mathcal{O})$ の $K$ に対し、
$i_g^*(Rf_{big, *}K) = Rf'_{small, *}(i_{g'}^*K)$
が $D(\textit{Mod}(X'_\etale, \mathcal{O}_{X'}))$ において成り立つ。
\item[(5)] $D((\textit{Spaces}/Y)_\etale)$ の $K$ に対し、\\
$g_{big}^{-1}(Rf_{big, *}K) = Rf'_{big, *}((g'_{big})^{-1}K)$
が $D((\textit{Spaces}/X')_\etale)$ において成り立つ。
\item[(6)] $D((\textit{Spaces}/Y)_\etale, \mathcal{O})$ の $K$ に対し、\\
$g_{big}^*(Rf_{big, *}K) = Rf'_{big, *}((g'_{big})^*K)$
が $D(\textit{Mod}(X'_\etale, \mathcal{O}_{X'}))$ において成り立つ。
\end{enumerate}
\end{lemma}

\begin{proof}
(1) は、$K$ を表すアーベル層の K-入射複体を選び、
補題 \ref{spaces-more-cohomology-lemma-compare-injectives} と
(\ref{spaces-more-cohomology-equation-compare-big-small}) を用いることから従う。

\medskip\noindent
(3) は、$K$ を表すアーベル層の K-入射複体を選び、
補題 \ref{spaces-more-cohomology-lemma-compare-injectives} と『トポロジー』補題
\ref{topologies-lemma-morphism-big-small-cartesian-diagram-etale}
を用いることから従う。

\medskip\noindent
(5) は『サイト上のコホモロジー』補題
\ref{sites-cohomology-lemma-localize-cartesian-square} である。

\medskip\noindent
(6) は『サイト上のコホモロジー』補題
\ref{sites-cohomology-lemma-localize-cartesian-square-modules} である。

\medskip\noindent
(2) は次のように証明できる。上で、環付きサイトの射として
$\pi_X \circ f_{big} = f_{small} \circ \pi_Y$ であることを見た。
したがって、『サイト上のコホモロジー』補題
\ref{sites-cohomology-lemma-derived-pushforward-composition} により
$R\pi_{X, *} \circ Rf_{big, *} = Rf_{small, *} \circ R\pi_{Y, *}$
を得る。制限関手 $\pi_{X, *}$ と $\pi_{Y, *}$ は完全なので結論を得る。

\medskip\noindent
(4) は、(6) と $f' : Y' \to X'$ に適用した (2) から従う。
\end{proof}

\noindent
$S$ をスキーム、$X$ を $S$ 上の代数空間とする。
$\mathcal{H}$ を $(\textit{Spaces}/X)_\etale$ 上のアーベル層とする。
$H^n_\etale(U, \mathcal{H})$ は $(\textit{Spaces}/X)_\etale$ の
対象 $U$ 上の $\mathcal{H}$ のコホモロジーを表すことを想起する。

\begin{lemma}
\label{spaces-more-cohomology-lemma-compare-cohomology}
$S$ をスキーム、$f : Y \to X$ を $S$ 上の代数空間の射とする。このとき
\begin{enumerate}
\item $D(X_\etale)$ の $K$ に対し
$H^n_\etale(X, \pi_X^{-1}K) = H^n(X_\etale, K)$ である。
\item $D(X_\etale, \mathcal{O}_X)$ の $K$ に対し
$H^n_\etale(X, L\pi_X^*K) = H^n(X_\etale, K)$ である。
\item $D(X_\etale)$ の $K$ に対し
$H^n_\etale(Y, \pi_X^{-1}K) = H^n(Y_\etale, f_{small}^{-1}K)$ である。
\item $D(X_\etale, \mathcal{O}_X)$ の $K$ に対し\\
$H^n_\etale(Y, L\pi_X^*K) = H^n(Y_\etale, Lf_{small}^*K)$ である。
\item $D((\textit{Spaces}/X)_\etale)$ の $M$ に対し
$H^n_\etale(Y, M) = H^n(Y_\etale, i_f^{-1}M)$ である。
\item $D((\textit{Spaces}/X)_\etale, \mathcal{O})$ の $M$ に対し
$H^n_\etale(Y, M) = H^n(Y_\etale, i_f^*M)$ である。
\end{enumerate}
\end{lemma}

\begin{proof}
(5) を証明するため、$M$ をアーベル層の K-入射複体で表し、
補題 \ref{spaces-more-cohomology-lemma-compare-injectives} を適用して定義を書き下す。
(3) は $i_f^{-1}\pi_X^{-1} = f_{small}^{-1}$ であることから従う。
(1) は (3) の特別な場合である。

\medskip\noindent
(6) は非常に一般的な『サイト上のコホモロジー』補題
\ref{sites-cohomology-lemma-pullback-same-cohomology} から従う。
次いで (4) は $Lf_{small}^* = i_f^* \circ L\pi_X^*$ から従う。
(2) は (4) の特別な場合である。
\end{proof}

\begin{lemma}
\label{spaces-more-cohomology-lemma-cohomological-descent-etale}
$S$ をスキーム、$X$ を $S$ 上の代数空間とする。
$K \in D(X_\etale)$ に対し、写像
$$
K \longrightarrow R\pi_{X, *}\pi_X^{-1}K
$$
は同型である。ここで
$\pi_X : \Sh((\textit{Spaces}/X)_\etale) \to \Sh(X_\etale)$ は上記の射である。
\end{lemma}

\begin{proof}
$\pi_X^{-1}$ と $\pi_{X, *} = i_X^{-1}$ はともに完全関手であり、
合成 $\pi_{X, *} \circ \pi_X^{-1}$ が恒等関手なので、主張は成り立つ。
\end{proof}

\begin{lemma}
\label{spaces-more-cohomology-lemma-compare-higher-direct-image-proper}
$S$ をスキームとする。$f : Y \to X$ を $S$ 上の代数空間の固有射とする。
このとき
\begin{enumerate}
\item 関手 $\Sh(Y_\etale) \to \Sh((\textit{Spaces}/X)_\etale)$ として
$\pi_X^{-1} \circ f_{small, *} = f_{big, *} \circ \pi_Y^{-1}$ であり、
\item コホモロジー層が捩れ層である $D^+(Y_\etale)$ の $K$ に対し\\
$\pi_X^{-1}Rf_{small, *}K = Rf_{big, *}\pi_Y^{-1}K$ であり、
\item $f$ が有限ならば、$D(Y_\etale)$ のすべての $K$ に対し
$\pi_X^{-1}Rf_{small, *}K = Rf_{big, *}\pi_Y^{-1}K$ である。
\end{enumerate}
\end{lemma}

\begin{proof}
(1) の証明。$\mathcal{F}$ を $Y_\etale$ 上の層とし、
$g : X' \to X$ を $(\textit{Spaces}/X)_\etale$ の対象とする。
ファイバー積
$$
\xymatrix{
Y' \ar[r]_{f'} \ar[d]_{g'} & X' \ar[d]^g \\
Y \ar[r]^f & X
}
$$
を考える。このとき
$$
\begin{aligned}
(f_{big, *}\pi_Y^{-1}\mathcal{F})(X')
&= (\pi_Y^{-1}\mathcal{F})(Y') \\
&= ((g'_{small})^{-1}\mathcal{F})(Y') \\
&= (f'_{small, *}(g'_{small})^{-1}\mathcal{F})(X')
\end{aligned}
$$
であり、第二の等号は補題 \ref{spaces-more-cohomology-lemma-describe-pullback} による。
一方、再び補題 \ref{spaces-more-cohomology-lemma-describe-pullback} により
$$
(\pi_X^{-1}f_{small, *}\mathcal{F})(X') =
(g_{small}^{-1}f_{small, *}\mathcal{F})(X')
$$
である。したがって、集合の層に対する固有基底変換
（補題 \ref{spaces-more-cohomology-lemma-proper-base-change-f-star}）により、
これら二つの集合は標準的に同型である。この同型は制限写像と両立し、同型
$\pi_X^{-1}f_{small, *}\mathcal{F} = f_{big, *}\pi_Y^{-1}\mathcal{F}$
を定める。よって関手の同型
$\pi_X^{-1} \circ f_{small, *} = f_{big, *} \circ \pi_Y^{-1}$ を得る。

\medskip\noindent
(2) の証明。$D(Y_\etale)$ の任意の $K$ に対し、標準基底変換写像\\
$\pi_X^{-1}Rf_{small, *}K \to Rf_{big, *}\pi_Y^{-1}K$ が存在する。
『サイト上のコホモロジー』注意
\ref{sites-cohomology-remark-base-change} を参照されたい。
これが同型であることを示すには、$(\Sch/X)_\etale$ の任意の対象
$g : X' \to X$ に対し、\\
$i_g : \Sh(X'_\etale) \to \Sh((\Sch/X)_\etale)$
による基底変換写像の逆像が同型であることを示せば十分である。
% P10-E0156
% P10-E0157
前段落と同じ $T', g', f'$ をとる。基底変換写像の逆像は
\begin{align*}
g_{small}^{-1}Rf_{small, *}K
& =
i_g^{-1}\pi_X^{-1}Rf_{small, *}K \\
& \to
i_g^{-1}Rf_{big, *}\pi_Y^{-1}K \\
& =
Rf'_{small, *}(i_{g'}^{-1}\pi_Y^{-1}K) \\
& =
Rf'_{small, *}((g'_{small})^{-1}K)
\end{align*}
である。ここで $\pi_X \circ i_g = g_{small}$、
$\pi_Y \circ i_{g'} = g'_{small}$、および補題
\ref{spaces-more-cohomology-lemma-compare-higher-direct-image} を用いた。
$K$ が下に有界で、その $K$ のコホモロジー層が捩れ層ならば、この写像は
固有基底変換定理（補題 \ref{spaces-more-cohomology-lemma-proper-base-change}）により同型である。

\medskip\noindent
(3) の証明。$f$ が有限ならば、関手 $f_{small, *}$ と $f_{big, *}$ は完全である。
$f_{small}$ については、これは『空間のコホモロジー』補題
\ref{spaces-cohomology-lemma-finite-higher-direct-image-zero} から従う。
$f$ の任意の基底変換 $f'$ も有限なので、補題
\ref{spaces-more-cohomology-lemma-compare-higher-direct-image} (3) から $f_{big, *}$ も完全である
（高次導来関手は零である）。したがって、この場合は (1) から従う。
\end{proof}



\section{fppf トポロジーとエタール・トポロジーの比較}
\label{spaces-more-cohomology-section-fppf-etale}

\noindent
本節は『エタール・コホモロジー』節
\ref{etale-cohomology-section-fppf-etale} の類似である。

\medskip\noindent
$S$ をスキーム、$X$ を $S$ 上の代数空間とする。
圏 $\textit{Spaces}/X$ 上で fppf トポロジーとエタール・トポロジーを考える。
恒等関手
$(\textit{Spaces}/X)_\etale \to (\textit{Spaces}/X)_{fppf}$
は連続であり、『サイト』命題 \ref{sites-proposition-get-morphism} を適用すると
サイトの射
$$
\epsilon_X :
(\textit{Spaces}/X)_{fppf} \longrightarrow (\textit{Spaces}/X)_\etale
$$
を定める。$\epsilon_{X, *}$ は台となる前層上では恒等関手であり、
$\epsilon_X^{-1}$ はエタール層にその fppf 層化を対応させることに注意する。
大エタール・サイトと小エタール・サイトを比較するサイトの射
$$
\pi_X : (\textit{Spaces}/X)_\etale \longrightarrow X_{spaces, \etale}
$$
を考える。節 \ref{spaces-more-cohomology-section-compare} を参照されたい。その合成はサイトの射
$$
a_X = \pi_X \circ \epsilon_X :
(\textit{Spaces}/X)_{fppf}
\longrightarrow
X_{spaces, \etale}
$$
を定める。$\mathcal{H}$ が $(\textit{Spaces}/X)_{fppf}$ 上のアーベル層ならば、
$(\textit{Spaces}/X)_{fppf}$ の対象 $U$ 上の $\mathcal{H}$ のコホモロジーを
$H^n_{fppf}(U, \mathcal{H})$ と書く。

\begin{lemma}
\label{spaces-more-cohomology-lemma-comparison-fppf-etale}
$S$ をスキーム、$X$ を $S$ 上の代数空間とする。
\begin{enumerate}
\item $\mathcal{F} \in \Sh(X_\etale)$ に対し
$\epsilon_{X, *}a_X^{-1}\mathcal{F} = \pi_X^{-1}\mathcal{F}$
かつ $a_{X, *}a_X^{-1}\mathcal{F} = \mathcal{F}$ である。
\item $\mathcal{F} \in \textit{Ab}(X_\etale)$ に対し、$i > 0$ ならば
$R^i\epsilon_{X, *}(a_X^{-1}\mathcal{F}) = 0$ である。
\end{enumerate}
\end{lemma}

\begin{proof}
$a_X^{-1}\mathcal{F} = \epsilon_X^{-1} \pi_X^{-1}\mathcal{F}$ である。
補題 \ref{spaces-more-cohomology-lemma-describe-pullback} により、エタール層
$\pi_X^{-1}\mathcal{F}$ は fppf トポロジーに対する層である。
したがって、これは $a_X^{-1}\mathcal{F}$ に等しい
（$\epsilon_X$ による逆像は fppf 層化で与えられる）。さらに、
$\epsilon_{X, *}$ は台となる前層上では恒等であることを想起する。
これで (1) は補題 \ref{spaces-more-cohomology-lemma-describe-pullback} における
$\pi_X^{-1}$ の明示的記述から直ちに従う。

\medskip\noindent
(2) はスキームの場合、すなわち『エタール・コホモロジー』補題
\ref{etale-cohomology-lemma-V-C-all-n-etale-fppf} (1) に帰着して証明する。
すべての代数空間はエタール局所的にスキームなので、これは「明らかにうまくいく」。
以下に詳細を記すが、読者には証明を飛ばすことを勧める。

\medskip\noindent
$(\textit{Spaces}/X)_{fppf}$ 上のアーベル層 $\mathcal{H}$ に対し、\\
高次順像 $R^p\epsilon_{X, *}\mathcal{H}$ は、
$(\textit{Spaces}/X)_\etale$ 上の前層\\
$U \mapsto H^p_{fppf}(U, \mathcal{H})$ に付随する層である。
『サイト上のコホモロジー』補題
\ref{sites-cohomology-lemma-higher-direct-images} を参照されたい。
$(\textit{Spaces}/X)_\etale$ のすべての対象はスキームによる被覆をもつので、
$U/X$ がスキームであり
$\xi \in H^p_{fppf}(U, a_X^{-1}\mathcal{F})$ が与えられたとき、
$\xi$ の $U_i$ への制限が零になるようなエタール被覆
$\{U_i \to U\}$ を見つけられることを示せば十分である。ここで
\begin{align*}
H^p_{fppf}(U, a_X^{-1}\mathcal{F})
& =
H^p((\textit{Spaces}/U)_{fppf}, (a_X^{-1}\mathcal{F})|_{\textit{Spaces}/U}) \\
& =
H^p((\Sch/U)_{fppf}, (a_X^{-1}\mathcal{F})|_{\Sch/U})
\end{align*}
である。第二の同一視は補題
\ref{spaces-more-cohomology-lemma-compare-cohomology-other-topologies} によるもので、第一は制限に関する一般論
（『サイト上のコホモロジー』補題
\ref{sites-cohomology-lemma-cohomology-of-open}）である。
最初の段落と、スキームの場合の対応する結果
（『エタール・コホモロジー』補題
\ref{etale-cohomology-lemma-describe-pullback-pi-fppf}）を見ると、層
$(a_X^{-1}\mathcal{F})|_{\Sch/U}$ は「$a_U$ のスキーム版」による逆像と一致する。
したがって、各 $i$ について我々の類が\\
$H^p((\Sch/U_i)_{fppf}, (a_X^{-1}\mathcal{F})|_{\Sch/U_i})$
で消えるようなエタール被覆 $\{U_i \to U\}$ を見つけられる。
『エタール・コホモロジー』補題
\ref{etale-cohomology-lemma-V-C-all-n-etale-fppf} を参照されたい
（ここで用いるべき正確な主張は、すべての $n$ について $V_n$ が成り立つことであり、
これはスキームの場合の (2) の主張である）。同じ公式を今度は $U_i$ に用いて
戻すと、望むとおり $\xi$ の $U_i$ 上への制限は零である。
\end{proof}

\noindent
スキームの場合に行われた難しい仕事から、小エタール・サイトに由来する層については
エタール・コホモロジーと fppf コホモロジーが一致することが分かる。

\begin{lemma}
\label{spaces-more-cohomology-lemma-cohomological-descent-etale-fppf}
$S$ をスキーム、$X$ を $S$ 上の代数空間とする。
$K \in D^+(X_\etale)$ に対し、写像
$$
\pi_X^{-1}K \longrightarrow R\epsilon_{X, *}a_X^{-1}K
\quad\text{および}\quad
K \longrightarrow Ra_{X, *}a_X^{-1}K
$$
は同型である。ここで
$a_X : \Sh((\textit{Spaces}/X)_{fppf}) \to \Sh(X_\etale)$ は上記の射である。
\end{lemma}

\begin{proof}
第二の主張だけを証明する。第一の主張はより容易で、全く同じ方法で証明される。
$K$ が単一のアーベル層で与えられる場合に直ちに帰着できる。
実際、$K$ を下に有界な複体 $\mathcal{F}^\bullet$ で表す。
層の場合から、$\mathcal{F}^n = a_{X, *} a_X^{-1} \mathcal{F}^n$ であり、
$q > 0$ に対して層 $R^qa_{X, *}a_X^{-1}\mathcal{F}^n$ は零である。
Leray の非輪状性補題
（『導来圏』補題 \ref{derived-lemma-leray-acyclicity}）を
$a_X^{-1}\mathcal{F}^\bullet$ と関手 $a_{X, *}$ に適用して結論を得る。
以下では $K = \mathcal{F}$ と仮定する。

\medskip\noindent
補題 \ref{spaces-more-cohomology-lemma-comparison-fppf-etale} により
$a_{X, *}a_X^{-1}\mathcal{F} = \mathcal{F}$ である。
したがって、$q > 0$ に対して
$R^qa_{X, *}a_X^{-1}\mathcal{F} = 0$ であることを示せば十分である。
このために
% P10-E0159
$a_X = \epsilon_X \circ \pi_X$ と Leray スペクトル系列
（『サイト上のコホモロジー』補題
\ref{sites-cohomology-lemma-relative-Leray}）を用いることができる。
補題 \ref{spaces-more-cohomology-lemma-comparison-fppf-etale} により、$i > 0$ に対して
$R^i\epsilon_{X, *}(a_X^{-1}\mathcal{F}) = 0$ である。
$\epsilon_{X, *}a_X^{-1}\mathcal{F} = \pi_X^{-1}\mathcal{F}$ であり、
補題 \ref{spaces-more-cohomology-lemma-cohomological-descent-etale} により、$j > 0$ に対して
$R^j\pi_{X, *}(\pi_X^{-1}\mathcal{F}) = 0$ である。これで証明が完了する。
\end{proof}

\begin{lemma}
\label{spaces-more-cohomology-lemma-compare-cohomology-etale-fppf}
$S$ をスキーム、$X$ を $S$ 上の代数空間とする。上記の
$a_X : \Sh((\textit{Spaces}/X)_{fppf}) \to \Sh(X_\etale)$ に対し、
\begin{enumerate}
\item $X_\etale$ 上のアーベル層 $\mathcal{F}$ について
$H^q(X_\etale, \mathcal{F}) = H^q_{fppf}(X, a_X^{-1}\mathcal{F})$ であり、
\item $K \in D^+(X_\etale)$ について
$H^q(X_\etale, K) = H^q_{fppf}(X, a_X^{-1}K)$ である。
\end{enumerate}
例：$A$ がアーベル群ならば
$H^q_\etale(X, \underline{A}) = H^q_{fppf}(X, \underline{A})$ である。
\end{lemma}

\begin{proof}
これは補題 \ref{spaces-more-cohomology-lemma-cohomological-descent-etale-fppf} と
『サイト上のコホモロジー』注意
\ref{sites-cohomology-remark-before-Leray} から従う。
\end{proof}

\begin{lemma}
\label{spaces-more-cohomology-lemma-push-pull-fppf-etale}
$S$ をスキームとし、$f : X \to Y$ を $S$ 上の代数空間の射とする。
このとき、トポスの可換図式
$$
\xymatrix{
\Sh((\textit{Spaces}/X)_{fppf}) \ar[rr]_{f_{big, fppf}} \ar[d]_{\epsilon_X} & &
\Sh((\textit{Spaces}/Y)_{fppf}) \ar[d]^{\epsilon_Y} \\
\Sh((\textit{Spaces}/X)_\etale) \ar[rr]^{f_{big, \etale}} & &
\Sh((\textit{Spaces}/Y)_\etale)
}
$$
および
$$
\xymatrix{
\Sh((\textit{Spaces}/X)_{fppf}) \ar[rr]_{f_{big, fppf}} \ar[d]_{a_X} & &
\Sh((\textit{Spaces}/Y)_{fppf}) \ar[d]^{a_Y} \\
\Sh(X_\etale) \ar[rr]^{f_{small}} & &
\Sh(Y_\etale)
}
$$
が存在し、
% P10-E0160
$a_X = \pi_X \circ \epsilon_X$ および $a_Y = \pi_X \circ \epsilon_X$ である。
\end{lemma}

\begin{proof}
これは関係する射の定義を書き下せば直ちに従う。
『空間上のトポロジー』節
\ref{spaces-topologies-section-fppf} および節 \ref{spaces-more-cohomology-section-compare} を参照されたい。
\end{proof}

\begin{lemma}
\label{spaces-more-cohomology-lemma-proper-push-pull-fppf-etale}
補題 \ref{spaces-more-cohomology-lemma-push-pull-fppf-etale} において $f$ が固有ならば、
\begin{enumerate}
\item $a_Y^{-1} \circ f_{small, *} = f_{big, fppf, *} \circ a_X^{-1}$ であり、
\item コホモロジー層が捩れ層である $D^+(X_\etale)$ の $K$ に対し\\
$a_Y^{-1}(Rf_{small, *}K) = Rf_{big, fppf, *}(a_X^{-1}K)$ である。
\end{enumerate}
\end{lemma}

\begin{proof}
(1) の証明。補題 \ref{spaces-more-cohomology-lemma-compare-higher-direct-image-proper} (1) の証明を
繰り返してもよいが、ここではその結果から導く。
$\epsilon_{Y, *}$ は台となる前層上では恒等関手なので、同型を反映する。
補題 \ref{spaces-more-cohomology-lemma-comparison-fppf-etale} により
$\epsilon_{Y, *} \circ a_Y^{-1} = \pi_Y^{-1}$ であり、$X$ についても同様である。
標準写像
$a_Y^{-1}f_{small, *}\mathcal{F} \to f_{big, fppf, *}a_X^{-1}\mathcal{F}$
が同型であることを示すには、
\begin{align*}
\pi_Y^{-1}f_{small, *}\mathcal{F}
& =
\epsilon_{Y, *}a_Y^{-1}f_{small, *}\mathcal{F} \\
& \to 
\epsilon_{Y, *}f_{big, fppf, *}a_X^{-1}\mathcal{F} \\
& =
f_{big, \etale, *} \epsilon_{X, *}a_X^{-1}\mathcal{F} \\
& =
f_{big, \etale, *}\pi_X^{-1}\mathcal{F}
\end{align*}
が同型であることを示せば十分である。これは補題
\ref{spaces-more-cohomology-lemma-compare-higher-direct-image-proper} (1) である。

\medskip\noindent
(2) を示すため、次を用いる。
\begin{align*}
R\epsilon_{Y, *}Rf_{big, fppf, *}a_X^{-1}K
& =
Rf_{big, \etale, *}R\epsilon_{X, *}a_X^{-1}K \\
& =
Rf_{big, \etale, *}\pi_X^{-1}K \\
& =
\pi_Y^{-1}Rf_{small, *}K \\
& =
R\epsilon_{Y, *} a_Y^{-1}Rf_{small, *}K
\end{align*}
% P10-E0158
% P10-E0161
第一の等号は補題 \ref{spaces-more-cohomology-lemma-push-pull-fppf-etale} の可換図式と
『サイト上のコホモロジー』補題
\ref{sites-cohomology-lemma-derived-pushforward-composition} による。
次の第二の等号は補題 \ref{spaces-more-cohomology-lemma-cohomological-descent-etale-fppf} である。
第三は補題 \ref{spaces-more-cohomology-lemma-compare-higher-direct-image-proper} (2) であり、
第四は再び補題 \ref{spaces-more-cohomology-lemma-cohomological-descent-etale-fppf} である。
したがって基底変換写像
$a_Y^{-1}(Rf_{small, *}K) \to Rf_{big, fppf, *}(a_X^{-1}K)$ は同型
$$
R\epsilon_{Y, *}a_Y^{-1}Rf_{small, *}K \to
R\epsilon_{Y, *}Rf_{big, fppf, *}a_X^{-1}K
$$
を誘導する。次の注意により証明が完了する：$L$ が $D^+(Y_\etale)$ に属し、
$M$ が $D^+((\textit{Spaces}/Y)_{fppf})$ に属する写像
$\alpha : a_Y^{-1}L \to M$ について、$R\epsilon_{Y, *}\alpha$ が同型ならば
この写像も同型である。実際、$H^i(\alpha)$ が同型であることを
$i$ について帰納的に示す。十分小さいすべての $i$ についてこれは成り立つ。
$i \leq i_0$ について成り立つとする。この範囲では
$H^i(M) = a_Y^{-1}H^i(L)$ なので、補題
\ref{spaces-more-cohomology-lemma-comparison-fppf-etale} により、$j > 0$ および $i \leq i_0$ に対し
$R^j\epsilon_{Y, *}H^i(M) = 0$ である。ゆえにスペクトル系列の議論により
$\epsilon_{Y, *}H^{i_0 + 1}(M) = H^{i_0 + 1}(R\epsilon_{Y, *}M)$ である。
したがって
$\epsilon_{Y, *}H^{i_0 + 1}(M) = \pi_Y^{-1}H^{i_0 + 1}(L) =
\epsilon_{Y, *}a_Y^{-1}H^{i_0 + 1}(L)$ である。
これは望むとおり $H^{i_0 + 1}(\alpha)$ が同型であることを意味する
（$\epsilon_{Y, *}$ は台となる前層上で恒等なので同型を反映する）。
\end{proof}

\begin{lemma}
\label{spaces-more-cohomology-lemma-finite-push-pull-fppf-etale}
補題 \ref{spaces-more-cohomology-lemma-push-pull-fppf-etale} において $f$ が有限ならば、
$D^+(X_\etale)$ の $K$ に対し
$a_Y^{-1}(Rf_{small, *}K) = Rf_{big, fppf, *}(a_X^{-1}K)$ である。
\end{lemma}

\begin{proof}
$V$ がスキームであるような全射エタール射 $V \to Y$ をとる。
基底変換写像が同型であることは、$V$ に制限した後で証明すれば十分である。
よって $Y$ がスキームであると仮定してよい。射は有限、したがって表現可能なので、
$X$ と $Y$ がともにスキームであると仮定してよい。この場合、結果は
スキームの場合（『エタール・コホモロジー』補題
\ref{etale-cohomology-lemma-V-C-all-n-etale-fppf} (2)）から、
節 \ref{spaces-more-cohomology-section-api}、特に補題
\ref{spaces-more-cohomology-lemma-compare-cohomology-other-topologies} で論じたトポスの比較を用いて従う。
いくつかの詳細は省略する。
\end{proof}

\begin{lemma}
\label{spaces-more-cohomology-lemma-descent-sheaf-fppf-etale}
補題 \ref{spaces-more-cohomology-lemma-push-pull-fppf-etale} において、$f$ が平坦、局所有限表示、
かつ全射であると仮定する。このとき関手
$$
\Sh(Y_\etale) \longrightarrow
\left\{
(\mathcal{G}, \mathcal{H}, \alpha)
\middle|
\begin{matrix}
\mathcal{G} \in \Sh(X_\etale),\ \mathcal{H} \in \Sh((\Sch/Y)_{fppf}), \\
\alpha : a_X^{-1}\mathcal{G} \to f_{big, fppf}^{-1}\mathcal{H}
\text{ は同型}
\end{matrix}
\right\}
$$
は、$\mathcal{F}$ を
$(f_{small}^{-1}\mathcal{F}, a_Y^{-1}\mathcal{F}, can)$ に送る同値である。
\end{lemma}

\begin{proof}
関手 $a_X^{-1}$ は完全忠実である（補題
\ref{spaces-more-cohomology-lemma-comparison-fppf-etale} により
$a_{X, *}a_X^{-1} = \text{id}$ である）。したがって忘却関手\\
$(\mathcal{G}, \mathcal{H}, \alpha) \mapsto \mathcal{H}$ は三つ組の圏を
$\Sh((\Sch/Y)_{fppf})$ の充満部分圏と同一視する。
さらに、関手 $a_Y^{-1}$ は完全忠実なので、補題の関手も完全忠実である。

\medskip\noindent
エタール被覆 $\{Y_i \to Y\}$ が与えられたとする。
$f_i : X_i \to Y_i$ を $f$ の基底変換とし、
$f_{ij} = f_i \times f_j : X_i \times_X X_j  \to Y_i \times_Y Y_j$ と記す。
主張：すべての $i, j$ について補題が $f_i$ と $f_{ij}$ に対して真ならば、
$f$ に対しても真である。これを見るため、与えられたエタール被覆が四つのサイト
$Y_\etale, X_\etale, (\Sch/Y)_{fppf}, (\Sch/X)_{fppf}$
のそれぞれにおける終対象のエタール被覆を定めることに注意する。
したがって、四つの場合のそれぞれで層の圏はこの被覆に関する貼り合わせデータの圏と
同値である（『サイト』補題 \ref{sites-lemma-mapping-property-glue}）。
ここで圏の巨大な可換図式が主張の証明を完了する。詳細は省略する。
この主張により $Y$ 上でエタール局所的に作業できる。特に、
$Y$ がスキームであると仮定してよい。

\medskip\noindent
$Y$ がスキームであると仮定する。スキーム $X'$ と全射エタール射
$s : X' \to X$ を選ぶ。$f' = f \circ s : X' \to Y$ とおくと、
$f'$ は全射、局所有限表示、かつ平坦である。
主張：補題が $f'$ に対して真ならば、$f$ に対しても真である。
実際、$f$ に対する三つ組 $(\mathcal{G}, \mathcal{H}, \alpha)$ が与えられたとき、
$s$ によって逆像をとると、$f'$ に対する三つ組
$(s_{small}^{-1}\mathcal{G}, \mathcal{H}, s_{big, fppf}^{-1}\alpha)$ を得る。
この三つ組に対する解は、$a_Y^{-1}\mathcal{F} = \mathcal{H}$ を満たす
$Y_\etale$ 上の層 $\mathcal{F}$ を与える。証明の最初の段落により、これは
三つ組が本質像に属することを意味する。これで $X$ と $Y$ がともに
スキームである場合に帰着される。この場合は『エタール・コホモロジー』補題
\ref{etale-cohomology-lemma-descent-sheaf-fppf-etale} から、
節 \ref{spaces-more-cohomology-section-api}、特に補題
\ref{spaces-more-cohomology-lemma-compare-cohomology-other-topologies} の議論を介して従う。
\end{proof}



\section{fppf トポロジーとエタール・トポロジーの比較：加群}
\label{spaces-more-cohomology-section-fppf-etale-modules}

\noindent
節 \ref{spaces-more-cohomology-section-fppf-etale} の議論を続けるが、本節では加群の層について
何が起こるかを簡潔に論じる。

\medskip\noindent
$S$ をスキーム、$X$ を $S$ 上の代数空間とする。節
\ref{spaces-more-cohomology-section-fppf-etale} で導入したサイトの射 $\epsilon_X$、$\pi_X$、
およびそれらの合成 $a_X$ は、環付きサイトの射へ自然に拡張される。
第一の射を
$$
\epsilon_X :
((\textit{Spaces}/X)_{fppf}, \mathcal{O})
\longrightarrow
((\textit{Spaces}/X)_\etale, \mathcal{O})
$$
と書く。実際、これらの層は同じ台となる前層をもつので、構造層に同じ記号を
用いてよいことに注意する。第二の射は
$$
\pi_X :
((\textit{Spaces}/X)_\etale, \mathcal{O})
\longrightarrow
(X_\etale, \mathcal{O}_X)
$$
である。第三の射は
$$
a_X :
((\textit{Spaces}/X)_{fppf}, \mathcal{O})
\longrightarrow
(X_\etale, \mathcal{O}_X)
$$
である。これらのサイト上の準連接加群について既知の事実を復習する。

\begin{lemma}
\label{spaces-more-cohomology-lemma-review-quasi-coherent}
$S$ をスキーム、$X$ を $S$ 上の代数空間とし、$\mathcal{F}$ を
準連接 $\mathcal{O}_X$-加群とする。
\begin{enumerate}
\item 規則
$$
\mathcal{F}^a : (\textit{Spaces}/X)_\etale \longrightarrow \textit{Ab},\quad
(f : Y \to X) \longmapsto \Gamma(Y, f^*\mathcal{F})
$$
は fpqc 被覆、したがって fppf 被覆およびエタール被覆に対する層条件を満たす。
\item $(\textit{Spaces}/X)_\etale$ 上で
$\mathcal{F}^a = \pi_X^*\mathcal{F}$ である。
\item $(\textit{Spaces}/X)_{fppf}$ 上で
$\mathcal{F}^a = a_X^*\mathcal{F}$ である。
\item 規則 $\mathcal{F} \mapsto \mathcal{F}^a$ は、準連接
$\mathcal{O}_X$-加群と
$((\textit{Spaces}/X)_\etale, \mathcal{O})$ 上の準連接加群との間の同値を定める。
\item 規則 $\mathcal{F} \mapsto \mathcal{F}^a$ は、準連接
$\mathcal{O}_X$-加群と
$((\textit{Spaces}/X)_{fppf}, \mathcal{O})$ 上の準連接加群との間の同値を定める。
\item $\epsilon_{X, *}a_X^*\mathcal{F} = \pi_X^*\mathcal{F}$ および
$a_{X, *}a_X^*\mathcal{F} = \mathcal{F}$ である。
\item $i > 0$ に対し $R^i\epsilon_{X, *}(a_X^*\mathcal{F}) = 0$ および
$R^ia_{X, *}(a_X^*\mathcal{F}) = 0$ である。
\end{enumerate}
\end{lemma}

\begin{proof}
(1) は準連接加群の fppf 降下の帰結である。実際、
$\{f_i : U_i \to U\}$ が $(\textit{Spaces}/X)_\etale$ における
fpqc 被覆であるとする。構造射を $g : U \to X$ と記す。
$s_i|_{U_i \times_U U_j} = s_j|_{U_i \times_U U_j}$ を満たす切断の族
$s_i \in \Gamma(U_i , f_i^*g^*\mathcal{F})$ があるとする。
% P10-E0163
対応する切断 $s \in \Gamma(U, g^*\mathcal{F})$ を見つけなければならない。
$s_i$ を写像の族
$\varphi_i : f_i^*\mathcal{O}_U = \mathcal{O}_{U_i} \to f_i^*g^*\mathcal{F}$
と解釈し直すことができ、これは $U$ 上の準連接層
$\mathcal{O}_U$ と $g^*\mathcal{F}$ に付随する標準降下データと両立する。
したがって、『空間上の降下』命題
\ref{spaces-descent-proposition-fpqc-descent-quasi-coherent} により、
これらを写像 $\mathcal{O}_U \to g^*\mathcal{F}$ に（一意に）降下でき、
それが求める切断 $s$ を与える。

\medskip\noindent
(2)--(7) はスキームに対する対応する主張から導く。
各 $X_i$ がスキームであるようなエタール被覆
$\{X_i \to X\}_{i \in I}$ を選ぶ。$X_i \times_X X_j$ もスキームである。
この被覆は三つのサイト
$(\textit{Spaces}/X)_{fppf}$、$(\textit{Spaces}/X)_\etale$、および $X_\etale$
のそれぞれで終対象の被覆を誘導する。したがって、これらのサイト上の層の圏は
これらの被覆に関する降下データと同値である。『サイト』補題
\ref{sites-lemma-mapping-property-glue} を参照されたい。
(2), (3) は局所的である（貼り合わせの主張があるため）。
準連接性は局所的性質なので、(4), (5) も局所的である。
明らかに (6), (7) も局所的である。したがって、$X$ がスキームである場合に
補題の (2)--(7) を証明すれば十分である。

\medskip\noindent
$X$ がスキームであると仮定する。埋め込み
$(\Sch/X)_\etale \subset (\textit{Spaces}/X)_\etale$ および
$(\Sch/X)_{fppf} \subset (\textit{Spaces}/X)_{fppf}$ は、補題
\ref{spaces-more-cohomology-lemma-compare-cohomology-other-topologies} により環付きトポスの同値を定める。
(2)--(7) はスキームの場合、すなわち『エタール・コホモロジー』補題
\ref{etale-cohomology-lemma-review-quasi-coherent} から従う。
この同値を介して準連接性を移すには、準連接性が加群の内在的性質であることを用いる。
『サイト上の加群』節 \ref{sites-modules-section-local} を参照されたい。
若干の詳細は省略する。
\end{proof}

\begin{lemma}
\label{spaces-more-cohomology-lemma-cohomological-descent-etale-fppf-modules}
$S$ をスキーム、$X$ を $S$ 上の代数空間とする。
準連接 $\mathcal{O}_X$-加群 $\mathcal{F}$ に対し、写像
$$
\pi_X^*\mathcal{F} \longrightarrow R\epsilon_{X, *}(a_X^*\mathcal{F})
\quad\text{および}\quad
\mathcal{F} \longrightarrow Ra_{X, *}(a_X^*\mathcal{F})
$$
は同型である。
\end{lemma}

\begin{proof}
これは補題 \ref{spaces-more-cohomology-lemma-review-quasi-coherent} (6), (7) の直ちに従う帰結である。
\end{proof}

\begin{lemma}
\label{spaces-more-cohomology-lemma-vanishing-adequate}
$S$ をスキーム、$X$ を $S$ 上の代数空間とする。
$\mathcal{F}_1 \to \mathcal{F}_2 \to \mathcal{F}_3$ を準連接
$\mathcal{O}_X$-加群の複体とする。$(\textit{Spaces}/X)_\etale$ 上で
$$
\mathcal{H}_\etale =
\Ker(\pi_X^*\mathcal{F}_2 \to \pi_X^*\mathcal{F}_3)/
\Im(\pi_X^*\mathcal{F}_1 \to \pi_X^*\mathcal{F}_2)
$$
とおき、$(\textit{Spaces}/X)_{fppf}$ 上で
$$
\mathcal{H}_{fppf} =
\Ker(a_X^*\mathcal{F}_2 \to a_X^*\mathcal{F}_3)/
\Im(a_X^*\mathcal{F}_1 \to a_X^*\mathcal{F}_2)
$$
とおく。このとき $\mathcal{H}_\etale = \epsilon_{X, *}\mathcal{H}_{fppf}$ であり、
$(\textit{Spaces}/X)_\etale$ の任意のアフィン対象 $U$ と $p > 0$ に対し
$$
H^p_\etale(U, \mathcal{H}_\etale) = H^p_{fppf}(U, \mathcal{H}_{fppf}) = 0
$$
である。
\end{lemma}

\noindent
より強いことが成り立つ。すなわち、$(\textit{Spaces}/X)_{fppf}$ 上の加群で、
fppf 局所的に補題の加群のように見えるものを
% P10-E0164
adequate 加群と呼ぶ。これらはすべての $\mathcal{O}$-加群の圏の弱 Serre 部分圏をなし、
そのコホモロジーは『Adequate 加群』節
\ref{adequate-section-adequate} で研究される。

\begin{proof}
$(\textit{Spaces}/X)_\etale$ の任意の対象 $f : U \to X$ に対し、
節 \ref{spaces-more-cohomology-section-compare} で論じた関手 $i_f^* = i_f^{-1}$ を介する
$\mathcal{H}_\etale$ の $U_\etale$ への制限
$\mathcal{H}_\etale|_{U_\etale}$ を考える。層
$\mathcal{H}_\etale|_{U_\etale}$ は複体 $f^*\mathcal{F}_\bullet$ の次数 $1$ の
ホモロジーに等しい。これは環付きサイト $U_\etale \to X_\etale$ の射として
$i_f \circ \pi_X = f$ であるためである。特に
$\mathcal{H}_\etale|_{U_\etale}$ は準連接 $\mathcal{O}_U$-加群である。
次に、$g : V \to U$ を $(\textit{Spaces}/X)_\etale$ における平坦射とする。
サイトの射 $V_\etale \to X_\etale$ として
$$
i_{f \circ g}^* \circ \pi_X^* = (f \circ g)^* = g^* \circ f^*
$$
であり、$g$ は平坦、したがって $g^*$ は完全なので、
$$
\mathcal{H}_\etale|_{V_\etale} =
g^*\left(\mathcal{H}_\etale|_{U_\etale}\right)
$$
を得る。
以上の準備により補題を証明できる。

\medskip\noindent
$f : U \to X$ を上記のものとし、
$\mathcal{U} = \{g_i : U_i \to U\}_{i \in I}$ を fppf 被覆とする。
補題 \ref{spaces-more-cohomology-lemma-review-quasi-coherent} (1) を
$\mathcal{H}_\etale|_{U_\etale}$ と上記の事実に適用すると、
$\mathcal{H}_\etale$ と被覆 $\mathcal{U}$ に対して層条件が成り立つ。
したがって $\mathcal{H}_\etale$ は既に fppf 層であり、これは前層として
$\mathcal{H}_{fppf}$ が $\mathcal{H}_\etale$ に等しいことを意味する。
特に $\mathcal{H}_\etale = \epsilon_{X, *}\mathcal{H}_{fppf}$ である。

\medskip\noindent
最後に消滅を証明するため、『サイト上のコホモロジー』補題
\ref{sites-cohomology-lemma-cech-vanish-collection} を用いる。
$\mathcal{B}$ を $(\textit{Spaces}/X)_{fppf}$ のアフィン対象全体とし、
$\text{Cov}$ を、$U$ と $U_i$ がアフィンである有限 fppf 被覆
$\mathcal{U} = \{U_i \to U\}_{i = 1, \ldots, n}$ 全体とする。このとき
$$
{\check H}^p(\mathcal{U}, \mathcal{H}_\etale) =
{\check H}^p(\mathcal{U}, \left(\mathcal{H}_\etale|_{U_\etale}\right)^a)
$$
である。実際、$U$ 上平坦なアフィン・スキーム
$U_{i_0} \times_U \ldots \times_U U_{i_p}$ 上での
$\mathcal{H}_\etale$ の値は、最初の段落により、準連接加群
$\mathcal{H}_\etale|_{U_\etale}$ の逆像の値と一致する。
したがって、『降下』補題
\ref{descent-lemma-standard-covering-Cech-quasi-coherent} により消滅を得る。
これで証明が完了する。
\end{proof}

\begin{lemma}
\label{spaces-more-cohomology-lemma-cohomological-descent-etale-fppf-modules-unbounded}
$S$ をスキーム、$X$ を $S$ 上の代数空間とする。
$K \in D_\QCoh(\mathcal{O}_X)$ に対し、写像
$$
L\pi_X^*K \longrightarrow R\epsilon_{X, *}(La_X^*K)
\quad\text{および}\quad
K \longrightarrow Ra_{X, *}(La_X^*K)
$$
は同型である。ここで
$a_X : \Sh((\textit{Spaces}/X)_{fppf}) \to \Sh(X_\etale)$ は上記の射である。
\end{lemma}

\begin{proof}
問題は $X$ 上エタール局所的なので、$X$ はアフィンであると仮定してよい。
$X = \Spec(A)$ とする。このとき『空間の導来圏』補題
\ref{spaces-perfect-lemma-derived-quasi-coherent-small-etale-site} および
『スキームの導来圏』補題 \ref{perfect-lemma-affine-compare-bounded} により
$D_\QCoh(\mathcal{O}_X) = D(A)$ である。したがって、$K$ を表す準連接
$\mathcal{O}_X$-加群の対応する複体 $\mathcal{K}^\bullet$ をもつ
$A$-加群の K-平坦複体
% P10-E0165
$K^\bullet$ を選べる。$\mathcal{K}^\bullet$ は
$\mathcal{O}_X$-加群の K-平坦複体であると主張する。

\medskip\noindent
主張の証明。『スキームの導来圏』補題
\ref{perfect-lemma-affine-K-flat} により、$\widetilde{K}^\bullet$ はスキーム
$(\Spec(A), \mathcal{O}_{\Spec(A)})$ 上 K-平坦である。
次に、$\epsilon$ を『空間の導来圏』補題
\ref{spaces-perfect-lemma-derived-quasi-coherent-small-etale-site} のものとすると
$\mathcal{K}^\bullet = \epsilon^*\widetilde{K}^\bullet$ である。
したがって、『サイト上のコホモロジー』補題
\ref{sites-cohomology-lemma-pullback-K-flat-points} と、スキームのエタール・サイトが
十分多くの点をもつこと（『エタール・コホモロジー』注意
\ref{etale-cohomology-remarks-enough-points}）により、
$\mathcal{K}^\bullet$ は K-平坦である。

\medskip\noindent
この主張により
$La_X^*K = a_X^*\mathcal{K}^\bullet$ および
$L\pi_X^*K = \pi_X^*\mathcal{K}^\bullet$ である。
証明の第一部は、準連接加群の逆像 $a_X^*\mathcal{K}^n$ が
$\epsilon_{X, *}$、それぞれ $a_{X, *}$ に対して非輪状であることを示している。
ならば Leray の非輪状性補題で証明は終わるはずではないか。
実はそうではない。Leray の非輪状性補題は下に有界な複体にしか適用できないからである。
しかし次の段落では、我々の複体が準連接加群の下に有界な複体の導来極限なので、
結果が下に有界な場合から従うことを示す。

\medskip\noindent
$\pi_X^*\mathcal{K}^\bullet$ と $a_X^*\mathcal{K}^\bullet$ のコホモロジー層は、
補題 \ref{spaces-more-cohomology-lemma-vanishing-adequate} により
$(\textit{Spaces}/X)_\etale$ のアフィン対象上で高次コホモロジー群が消滅する。
したがって、『サイト上のコホモロジー』補題
\ref{sites-cohomology-lemma-is-limit-dimension} により
$$
L\pi_X^*K = R\lim \tau_{\geq -n}(L\pi_X^*K)
\quad\text{および}\quad
La_X^*K = R\lim \tau_{\geq -n}(La_X^*K)
$$
である。

\medskip\noindent
% P10-E0166
$L\pi_X^*K = R\epsilon_{X, *}(La_X^*\mathcal{F})$ の証明。
上記により
$$
R\epsilon_{X, *}La_X^*K =
R\lim R\epsilon_{X, *}(\tau_{\geq -n}(La_X^*K))
$$
である。『サイト上のコホモロジー』補題
\ref{sites-cohomology-lemma-Rf-commutes-with-Rlim} を参照されたい。
$\tau_{\geq -n}(La_X^*K)$ は
$\tau_{\geq -n}(a_X^*\mathcal{K}^\bullet)$ で表されるが、これは
$a_X^*(\tau_{\geq -n}\mathcal{K}^\bullet)$ と同じとは限らない。
しかし、系
$$
\{\tau_{\geq -n}(a_X^*\mathcal{K}^\bullet)\}_{n \geq 1}
\quad\text{および}\quad
\{a_X^*(\tau_{\geq -n}\mathcal{K}^\bullet)\}_{n \geq 1}
$$
は pro-系として同型であることは明らかである。
Leray の非輪状性補題
（『導来圏』補題 \ref{derived-lemma-leray-acyclicity}）と補題の第一部により
$$
R\epsilon_{X, *}(a_X^*(\tau_{\geq -n}\mathcal{K}^\bullet)) =
\pi_X^*(\tau_{\geq -n}\mathcal{K}^\bullet)
$$
である。次に、系
$$
\{\tau_{\geq -n}(\pi_X^*\mathcal{K}^\bullet)\}_{n \geq 1}
\quad\text{および}\quad
\{\pi_X^*(\tau_{\geq -n}\mathcal{K}^\bullet)\}_{n \geq 1}
$$
が pro-系として同型であることを用いる。最後に、すべてを次のようにまとめる。
\begin{align*}
R\epsilon_{X, *}La_X^*K
& =
R\epsilon_{X, *} (R\lim \tau_{\geq -n}(La_X^*K)) \\
& =
R\lim R\epsilon_{X, *}(\tau_{\geq -n}(La_X^*K)) \\
& =
R\lim R\epsilon_{X, *}(\tau_{\geq -n}(a_X^*\mathcal{K}^\bullet)) \\
& =
R\lim R\epsilon_{X, *}(a_X^*(\tau_{\geq -n}\mathcal{K}^\bullet)) \\
& =
R\lim \pi_X^*(\tau_{\geq -n}\mathcal{K}^\bullet) \\
& =
R\lim \tau_{\geq -n}(\pi_X^*\mathcal{K}^\bullet) \\
& =
R\lim \tau_{\geq -n}(L\pi_X^*K) \\
& =
L\pi_X^*K
\end{align*}
ここで第四および第六の等号では、同型な pro-系は同じ $R\lim$ をもつことを用いた
（細部は省略する）。補題 \ref{spaces-more-cohomology-lemma-vanishing-adequate} で証明された複体
$\tau_{\geq -n}a_X^*\mathcal{K}^\bullet$ の項のコホモロジーについて
より多くを用いれば、この段階を避けられる。実際、そこから直接
$R\epsilon_{X, *}(\tau_{\geq -n}(a_X^*\mathcal{K}^\bullet)) =
\tau_{\geq -n}(\pi_X^*\mathcal{K}^\bullet)$ が従う。

\medskip\noindent
% P10-E0167
等式 $K = Ra_{X, *}(La_X^*\mathcal{F})$ も全く同じ方法で証明される。
最後の段階で『空間の導来圏』補題
\ref{spaces-perfect-lemma-nice-K-injective} による
$K = R\lim \tau_{\geq -n}K$ を用いる。
\end{proof}







\section{ph トポロジーとエタール・トポロジーの比較}
\label{spaces-more-cohomology-section-ph-etale}

\noindent
本節は『エタール・コホモロジー』節
\ref{etale-cohomology-section-ph-etale} の類似である。

\medskip\noindent
$S$ をスキーム、$X$ を $S$ 上の代数空間とする。
圏 $\textit{Spaces}/X$ 上で ph トポロジーとエタール・トポロジーを考える。
『空間上のトポロジー』補題
\ref{spaces-topologies-lemma-zariski-etale-smooth-syntomic-fppf-ph} により、
すべてのエタール被覆は ph 被覆なので、恒等関手
$(\textit{Spaces}/X)_\etale \to (\textit{Spaces}/X)_{ph}$ は連続である。
したがって、『サイト』命題 \ref{sites-proposition-get-morphism} を適用すると、
サイトの射
$$
\epsilon_X :
(\textit{Spaces}/X)_{ph} \longrightarrow (\textit{Spaces}/X)_\etale
$$
を定める。$\epsilon_{X, *}$ は台となる前層上では恒等関手であり、
$\epsilon_X^{-1}$ はエタール層にその ph 層化を対応させることに注意する。
大エタール・サイトと小エタール・サイトを比較するサイトの射
$$
\pi_X : (\textit{Spaces}/X)_\etale \longrightarrow X_{spaces, \etale}
$$
を考える。節 \ref{spaces-more-cohomology-section-compare} を参照されたい。その合成はサイトの射
$$
a_X = \pi_X \circ \epsilon_X :
(\textit{Spaces}/X)_{ph}
\longrightarrow
X_{spaces, \etale}
$$
を定める。$\mathcal{H}$ が $(\textit{Spaces}/X)_{ph}$ 上のアーベル層ならば、
$(\textit{Spaces}/X)_{ph}$ の対象 $U$ 上の $\mathcal{H}$ のコホモロジーを
$H^n_{ph}(U, \mathcal{H})$ と書く。

\begin{lemma}
\label{spaces-more-cohomology-lemma-comparison-ph-etale}
$S$ をスキーム、$X$ を $S$ 上の代数空間とする。
\begin{enumerate}
\item $\mathcal{F} \in \Sh(X_\etale)$ に対し
$\epsilon_{X, *}a_X^{-1}\mathcal{F} = \pi_X^{-1}\mathcal{F}$
かつ $a_{X, *}a_X^{-1}\mathcal{F} = \mathcal{F}$ である。
\item 捩れ層 $\mathcal{F} \in \textit{Ab}(X_\etale)$ に対し、$i > 0$ ならば
$R^i\epsilon_{X, *}(a_X^{-1}\mathcal{F}) = 0$ である。
\end{enumerate}
\end{lemma}

\begin{proof}
$a_X^{-1}\mathcal{F} = \epsilon_X^{-1} \pi_X^{-1}\mathcal{F}$ である。
補題 \ref{spaces-more-cohomology-lemma-describe-pullback} により、エタール層
$\pi_X^{-1}\mathcal{F}$ は ph トポロジーに対する層である。
したがって、これは $a_X^{-1}\mathcal{F}$ に等しい
（$\epsilon_X$ による逆像は ph 層化で与えられる）。さらに、
$\epsilon_{X, *}$ は台となる前層上では恒等であることを想起する。
これで (1) は補題 \ref{spaces-more-cohomology-lemma-describe-pullback} における
$\pi_X^{-1}$ の明示的記述から直ちに従う。

\medskip\noindent
(2) はスキームの場合、すなわち『エタール・コホモロジー』補題
\ref{etale-cohomology-lemma-V-C-all-n-etale-ph} (1) に帰着して証明する。
すべての代数空間はエタール局所的にスキームなので、これは「明らかにうまくいく」。
以下に詳細を記すが、読者には証明を飛ばすことを勧める。

\medskip\noindent
$(\textit{Spaces}/X)_{ph}$ 上のアーベル層 $\mathcal{H}$ に対し、\\
高次順像 $R^p\epsilon_{X, *}\mathcal{H}$ は
$(\textit{Spaces}/X)_\etale$ 上の前層\\
$U \mapsto H^p_{ph}(U, \mathcal{H})$ に付随する層である。
『サイト上のコホモロジー』補題
\ref{sites-cohomology-lemma-higher-direct-images} を参照されたい。
$(\textit{Spaces}/X)_\etale$ のすべての対象はスキームによる被覆をもつので、
$U/X$ がスキームであり
$\xi \in H^p_{ph}(U, a_X^{-1}\mathcal{F})$ が与えられたとき、
$\xi$ の $U_i$ への制限が零になるようなエタール被覆
$\{U_i \to U\}$ を見つけられることを示せば十分である。ここで
\begin{align*}
H^p_{ph}(U, a_X^{-1}\mathcal{F})
& =
H^p((\textit{Spaces}/U)_{ph}, (a_X^{-1}\mathcal{F})|_{\textit{Spaces}/U}) \\
& =
H^p((\Sch/U)_{ph}, (a_X^{-1}\mathcal{F})|_{\Sch/U})
\end{align*}
である。第二の同一視は補題
\ref{spaces-more-cohomology-lemma-compare-cohomology-other-topologies} によるもので、第一は制限に関する一般論
（『サイト上のコホモロジー』補題
\ref{sites-cohomology-lemma-cohomology-of-open}）である。
最初の段落と、スキームの場合の対応する結果
（『エタール・コホモロジー』補題
\ref{etale-cohomology-lemma-describe-pullback-pi-ph}）を見ると、層
$(a_X^{-1}\mathcal{F})|_{\Sch/U}$ は「$a_U$ のスキーム版」による逆像と一致する。
したがって、各 $i$ について我々の類が
$H^p((\Sch/U_i)_{ph}, (a_X^{-1}\mathcal{F})|_{\Sch/U_i})$
で消えるようなエタール被覆 $\{U_i \to U\}$ を見つけられる。
『エタール・コホモロジー』補題
\ref{etale-cohomology-lemma-V-C-all-n-etale-ph} を参照されたい
（ここで用いるべき正確な主張は、すべての $n$ について $V_n$ が成り立つことであり、
これはスキームの場合の (2) の主張である）。同じ公式を今度は $U_i$ に用いて
戻すと、望むとおり $\xi$ の $U_i$ 上への制限は零である。
\end{proof}

\noindent
スキームの場合に行われた難しい仕事から、小エタール・サイトに由来する
アーベル捩れ層については、エタール・コホモロジーと ph コホモロジーが一致する。

\begin{lemma}
\label{spaces-more-cohomology-lemma-cohomological-descent-etale-ph}
$S$ をスキーム、$X$ を $S$ 上の代数空間とする。
コホモロジー層が捩れ層である $K \in D^+(X_\etale)$ に対し、写像
$$
\pi_X^{-1}K \longrightarrow R\epsilon_{X, *}a_X^{-1}K
\quad\text{および}\quad
K \longrightarrow Ra_{X, *}a_X^{-1}K
$$
は同型である。ここで
$a_X : \Sh((\textit{Spaces}/X)_{ph}) \to \Sh(X_\etale)$ は上記の射である。
\end{lemma}

\begin{proof}
第二の主張だけを証明する。第一の主張はより容易で、全く同じ方法で証明される。
$K$ が単一のアーベル捩れ層で与えられる場合に帰着できる。
実際、$K$ をアーベル捩れ層の下に有界な複体 $\mathcal{F}^\bullet$ で表す。
これは『サイト上のコホモロジー』補題
\ref{sites-cohomology-lemma-torsion} により可能である。
層の場合から、$\mathcal{F}^n = a_{X, *} a_X^{-1} \mathcal{F}^n$ であり、
$q > 0$ に対して層 $R^qa_{X, *}a_X^{-1}\mathcal{F}^n$ は零である。
Leray の非輪状性補題
（『導来圏』補題 \ref{derived-lemma-leray-acyclicity}）を
$a_X^{-1}\mathcal{F}^\bullet$ と関手 $a_{X, *}$ に適用して結論を得る。
以下では $K = \mathcal{F}$ とし、$\mathcal{F}$ はアーベル捩れ層であると仮定する。

\medskip\noindent
補題 \ref{spaces-more-cohomology-lemma-comparison-ph-etale} により
$a_{X, *}a_X^{-1}\mathcal{F} = \mathcal{F}$ である。
したがって、$q > 0$ に対し
$R^qa_{X, *}a_X^{-1}\mathcal{F} = 0$ であることを示せば十分である。
このために
% P10-E0168
$a_X = \epsilon_X \circ \pi_X$ と Leray スペクトル系列
（『サイト上のコホモロジー』補題
\ref{sites-cohomology-lemma-relative-Leray}）を用いることができる。
補題 \ref{spaces-more-cohomology-lemma-comparison-ph-etale} により、$i > 0$ に対し
$R^i\epsilon_{X, *}(a_X^{-1}\mathcal{F}) = 0$ である。
$\epsilon_{X, *}a_X^{-1}\mathcal{F} = \pi_X^{-1}\mathcal{F}$ であり、
補題 \ref{spaces-more-cohomology-lemma-cohomological-descent-etale} により、$j > 0$ に対し
$R^j\pi_{X, *}(\pi_X^{-1}\mathcal{F}) = 0$ である。これで証明が完了する。
\end{proof}
\begin{lemma}
\label{spaces-more-cohomology-lemma-compare-cohomology-etale-ph}
$S$ をスキーム、$X$ を $S$ 上の代数空間とする。上記の
$a_X : \Sh((\textit{Spaces}/X)_{ph}) \to \Sh(X_\etale)$ に対し、
\begin{enumerate}
\item $X_\etale$ 上のアーベル捩れ層 $\mathcal{F}$ について\\
$H^q(X_\etale, \mathcal{F}) = H^q_{ph}(X, a_X^{-1}\mathcal{F})$ であり、
\item コホモロジー層が捩れ層である $K \in D^+(X_\etale)$ について\\
$H^q(X_\etale, K) = H^q_{ph}(X, a_X^{-1}K)$ である。
\end{enumerate}
例：$A$ がアーベル捩れ群ならば
$H^q_\etale(X, \underline{A}) = H^q_{ph}(X, \underline{A})$ である。
\end{lemma}

\begin{proof}
これは補題 \ref{spaces-more-cohomology-lemma-cohomological-descent-etale-ph} と
『サイト上のコホモロジー』注意
\ref{sites-cohomology-remark-before-Leray} から従う。
\end{proof}

\begin{lemma}
\label{spaces-more-cohomology-lemma-push-pull-ph-etale}
$S$ をスキームとし、$f : X \to Y$ を $S$ 上の代数空間の射とする。
このとき、トポスの可換図式
$$
\xymatrix{
\Sh((\textit{Spaces}/X)_{ph}) \ar[rr]_{f_{big, ph}} \ar[d]_{\epsilon_X} & &
\Sh((\textit{Spaces}/Y)_{ph}) \ar[d]^{\epsilon_Y} \\
\Sh((\textit{Spaces}/X)_\etale) \ar[rr]^{f_{big, \etale}} & &
\Sh((\textit{Spaces}/Y)_\etale)
}
$$
および
$$
\xymatrix{
\Sh((\textit{Spaces}/X)_{ph}) \ar[rr]_{f_{big, ph}} \ar[d]_{a_X} & &
\Sh((\textit{Spaces}/Y)_{ph}) \ar[d]^{a_Y} \\
\Sh(X_\etale) \ar[rr]^{f_{small}} & &
\Sh(Y_\etale)
}
$$
が存在し、
% P10-E0169
$a_X = \pi_X \circ \epsilon_X$ および $a_Y = \pi_X \circ \epsilon_X$ である。
\end{lemma}

\begin{proof}
これは関係する射の定義を書き下せば直ちに従う。
『空間上のトポロジー』節
\ref{spaces-topologies-section-ph} および節 \ref{spaces-more-cohomology-section-compare} を参照されたい。
\end{proof}

\begin{lemma}
\label{spaces-more-cohomology-lemma-proper-push-pull-ph-etale}
補題 \ref{spaces-more-cohomology-lemma-push-pull-ph-etale} において $f$ が固有ならば、
\begin{enumerate}
\item $a_Y^{-1} \circ f_{small, *} = f_{big, ph, *} \circ a_X^{-1}$ であり、
\item コホモロジー層が捩れ層である $D^+(X_\etale)$ の $K$ に対し\\
$a_Y^{-1}(Rf_{small, *}K) = Rf_{big, ph, *}(a_X^{-1}K)$ である。
\end{enumerate}
\end{lemma}

\begin{proof}
(1) の証明。補題 \ref{spaces-more-cohomology-lemma-compare-higher-direct-image-proper} (1) の証明を
繰り返してもよいが、ここではその結果から導く。
$\epsilon_{Y, *}$ は台となる前層上では恒等関手なので、同型を反映する。
補題 \ref{spaces-more-cohomology-lemma-comparison-ph-etale} により
$\epsilon_{Y, *} \circ a_Y^{-1} = \pi_Y^{-1}$ であり、$X$ についても同様である。
標準写像
$a_Y^{-1}f_{small, *}\mathcal{F} \to f_{big, ph, *}a_X^{-1}\mathcal{F}$
が同型であることを示すには、
\begin{align*}
\pi_Y^{-1}f_{small, *}\mathcal{F}
& =
\epsilon_{Y, *}a_Y^{-1}f_{small, *}\mathcal{F} \\
& \to 
\epsilon_{Y, *}f_{big, ph, *}a_X^{-1}\mathcal{F} \\
& =
f_{big, \etale, *} \epsilon_{X, *}a_X^{-1}\mathcal{F} \\
& =
f_{big, \etale, *}\pi_X^{-1}\mathcal{F}
\end{align*}
が同型であることを示せば十分である。これは補題
\ref{spaces-more-cohomology-lemma-compare-higher-direct-image-proper} (1) である。

\medskip\noindent
(2) を示すため、次を用いる。
\begin{align*}
R\epsilon_{Y, *}Rf_{big, ph, *}a_X^{-1}K
& =
Rf_{big, \etale, *}R\epsilon_{X, *}a_X^{-1}K \\
& =
Rf_{big, \etale, *}\pi_X^{-1}K \\
& =
\pi_Y^{-1}Rf_{small, *}K \\
& =
R\epsilon_{Y, *} a_Y^{-1}Rf_{small, *}K
\end{align*}
% P10-E0170
% P10-E0171
第一の等号は補題 \ref{spaces-more-cohomology-lemma-push-pull-ph-etale} の可換図式と
『サイト上のコホモロジー』補題
\ref{sites-cohomology-lemma-derived-pushforward-composition} による。
次の第二の等号は補題 \ref{spaces-more-cohomology-lemma-cohomological-descent-etale-ph} である。
第三は補題 \ref{spaces-more-cohomology-lemma-compare-higher-direct-image-proper} (2) であり、
第四は再び補題 \ref{spaces-more-cohomology-lemma-cohomological-descent-etale-ph} である。
したがって基底変換写像
$a_Y^{-1}(Rf_{small, *}K) \to Rf_{big, ph, *}(a_X^{-1}K)$ は同型
$$
R\epsilon_{Y, *}a_Y^{-1}Rf_{small, *}K \to
R\epsilon_{Y, *}Rf_{big, ph, *}a_X^{-1}K
$$
を誘導する。次の注意により証明が完了する：コホモロジー層が捩れ層である
$D^+(Y_\etale)$ の $L$ と、$D^+((\textit{Spaces}/Y)_{ph})$ の $M$ をもち、
$\alpha : a_Y^{-1}L \to M$ である写像を考える。
$R\epsilon_{Y, *}\alpha$ が同型ならば、$\alpha$ は同型である。
実際、$H^i(\alpha)$ が同型であることを $i$ について帰納的に示す。
十分小さいすべての $i$ についてこれは成り立つ。
$i \leq i_0$ について成り立つとする。この範囲では
$H^i(M) = a_Y^{-1}H^i(L)$ なので、補題
\ref{spaces-more-cohomology-lemma-comparison-ph-etale} により、$j > 0$ および $i \leq i_0$ に対し
$R^j\epsilon_{Y, *}H^i(M) = 0$ である。ゆえにスペクトル系列の議論により
$\epsilon_{Y, *}H^{i_0 + 1}(M) = H^{i_0 + 1}(R\epsilon_{Y, *}M)$ である。
したがって
$\epsilon_{Y, *}H^{i_0 + 1}(M) = \pi_Y^{-1}H^{i_0 + 1}(L) =
\epsilon_{Y, *}a_Y^{-1}H^{i_0 + 1}(L)$ である。
これは望むとおり $H^{i_0 + 1}(\alpha)$ が同型であることを意味する
（$\epsilon_{Y, *}$ は台となる前層上で恒等なので同型を反映する）。
\end{proof}
